% English translation of checked_pages/esq_119_011.tex; source: cote 119, site 11, PDF page 12, typescript page 9.
% AI-translated by GPT-6 Astra (OpenAI, ChatGPT Pro mode).
\scanpage{119/11}
\mspage{9}

allow these surfaces to be considered as associated with
canonically determined complex algebraic curves (and even curves over $\overline{\mathbb{Q}}$).
Here there is a typical interplay between the combinatorial
and the complex algebraic (or better, the algebraic over $\overline{\mathbb{Q}}$).

The ``small box'', with its identical pieces, which has
the charm of economy, will probably give a relatively
complicated description of the relations (complicated, but by no means
inextricable !). The large box will give rise to more numerous
relations (because there are many more base points and
remarkable paths between them), but with a more transparent structure. I
expect that in moduli dimension 2, just as in the more or
less familiar case of moduli dimension 1 (with, in particular, the
description of $S\ell(2,\mathbb{Z})$ by $(\rho,\sigma \mid \rho^3 = \sigma^2,\ \rho^4 = \sigma^6 = 1)$), one will find
generation by the automorphism groups of the three relevant types
of pieces, with simple relations that I have not worked out
at the time of writing these lines. Perhaps one will even find a
principle of this kind for all the $T_{g,\nu}$, as well as a cellular
decomposition of $M^{\wedge}_{g,\nu}$ generalizing those that present themselves spontaneously
for $M^{\wedge}_{0,4}$ and $M^{\wedge}_{1,1}$, and that I already glimpse for
moduli dimension 2, using the hypersurfaces corresponding
to the various \uline{real structures} on the complex structures
under consideration to carry out the desired cellular decomposition.


3. Instead of following (as planned) a rigorous thematic
order, I have let myself be carried away by my predilection for a
particularly rich and burning theme, to which I moreover intend to
devote myself as a priority for some time, beginning with the start
of the 84/85 academic year. I therefore resume the thematic account where I left it,
at the very beginning of the preceding paragraph.

My interest in topological surfaces begins to emerge
in 1974, when I propose to Yves Ladegaillerie the theme of studying
isotopically the embeddings of a topological 1-complex in a compact
surface. In the two years that follow, this study leads him to
a remarkable isotopy theorem, giving a complete algebraic description

\margin{11}

